\documentclass[11pt]{article}
\usepackage[letterpaper,margin=1in]{geometry}
\usepackage{amsmath,amssymb,amsthm}
\usepackage{booktabs,tabularx,array}
\usepackage[colorlinks=true,linkcolor=black,citecolor=black,urlcolor=blue]{hyperref}
\hypersetup{
  pdftitle={An acyclic edge coloring of K8 with no acyclic completion},
  pdfauthor={Alex Chan}
}
\newtheorem{theorem}{Theorem}
\newtheorem{lemma}{Lemma}
\newcommand{\out}{N^{+}}

\title{An acyclic edge coloring of \(K_8\)\\
with no acyclic completion}
\author{Alex Chan\\[2pt]
  \small Harvard University and National Bureau of Economic Research\\[-1pt]
  \small \texttt{achan@hbs.edu}}
\date{\today}

\begin{document}
\maketitle

\begin{abstract}
Alon, Sudakov, and Zaks suggested that every acyclic edge coloring
of \(K_{2n}\) with \(2n+1\) colors might extend to an acyclic edge
coloring of \(K_{2n+1}\) without introducing new colors.
We give an explicit counterexample on \(K_8\).
The coloring has exactly two proper extensions to \(K_9\) with the
same palette, and each contains a cycle using only colors \(2\) and \(4\).
We provide certificates for all two-color subgraphs, so the proof
can be checked directly from the displayed matrix.
\end{abstract}

\section{Introduction}

An edge coloring of a finite simple graph is \emph{acyclic} if it
is proper and the union of every two color classes is a forest.
In Section~1.2, printed page~160, Alon, Sudakov, and Zaks~\cite{ASZ}
suggested that every acyclic edge coloring of \(K_{2n}\) with
\(2n+1\) colors might be completed to an acyclic edge coloring of
\(K_{2n+1}\) without introducing new colors.
Here completion preserves the colors of all existing edges.

We give an explicit coloring for which such a completion is
impossible.

\begin{theorem}\label{thm:main}
There is a proper acyclic edge coloring of \(K_8\) with palette
\(\{1,\ldots,9\}\) that has no proper acyclic extension to \(K_9\)
with the same palette.
\end{theorem}

The construction and the nonextension statement have also been
formalized in Lean.\footnote{The source files are available at
\url{https://github.com/alexgocardinal/acyclic-completion}.
The theorem
\texttt{AcyclicCompletion.certified\_counterexample}
proves properness, acyclicity, and nonextension, using the project's
explicit encodings of edge colorings and finite cycles.
The formalization uses Lean's standard library.}

\section{The coloring}

Let the vertices of \(K_8\) be \(0,\ldots,7\).
For distinct \(u,v\), define \(c(uv)\) by the following symmetric
matrix; the diagonal entries do not represent edges.
\[
\begin{array}{c|rrrrrrrr}
 &0&1&2&3&4&5&6&7\\ \hline
0&-&7&8&6&4&3&1&5\\
1&7&-&6&1&8&9&5&4\\
2&8&6&-&9&1&5&2&7\\
3&6&1&9&-&7&8&3&2\\
4&4&8&1&7&-&2&9&3\\
5&3&9&5&8&2&-&4&1\\
6&1&5&2&3&9&4&-&6\\
7&5&4&7&2&3&1&6&-
\end{array}
\]

\begin{lemma}\label{lem:acyclic}
The coloring \(c\) is proper and acyclic.
\end{lemma}

\begin{proof}
The seven entries off the diagonal in each row are distinct,
so \(c\) is proper.
Table~\ref{tab:paths} lists the path components of every two-color
subgraph, including isolated vertices.
For each pair, the listed paths have disjoint vertex sets covering
all eight vertices, and their consecutive edges are exactly the
edges carrying either of the two colors in the matrix.
Thus each two-color subgraph is a disjoint union of paths and
is a forest.
\end{proof}

\begin{table}[htbp]
\centering
\small
\setlength{\tabcolsep}{8pt}
\begin{tabular}{clcl}
\toprule
Colors & Path components & Colors & Path components\\
\midrule
\(1,2\)&\texttt{06245731}&\(3,7\)&\texttt{105; 27436}\\
\(1,3\)&\texttt{13605742}&\(3,8\)&\texttt{147; 20536}\\
\(1,4\)&\texttt{24065713}&\(3,9\)&\texttt{051; 23647}\\
\(1,5\)&\texttt{31607524}&\(4,5\)&\texttt{2561704; 3}\\
\(1,6\)&\texttt{42130675}&\(4,6\)&\texttt{21765; 304}\\
\(1,7\)&\texttt{57243106}&\(4,7\)&\texttt{271043; 56}\\
\(1,8\)&\texttt{60241357}&\(4,8\)&\texttt{20417; 356}\\
\(1,9\)&\texttt{06423157}&\(4,9\)&\texttt{046517; 23}\\
\(2,3\)&\texttt{0547362; 1}&\(5,6\)&\texttt{3076125; 4}\\
\(2,4\)&\texttt{04562; 173}&\(5,7\)&\texttt{34; 527016}\\
\(2,5\)&\texttt{073; 16254}&\(5,8\)&\texttt{35207; 416}\\
\(2,6\)&\texttt{037621; 45}&\(5,9\)&\texttt{07; 325164}\\
\(2,7\)&\texttt{01; 543726}&\(6,7\)&\texttt{4301276; 5}\\
\(2,8\)&\texttt{026; 14537}&\(6,8\)&\texttt{412035; 67}\\
\(2,9\)&\texttt{0; 1546237}&\(6,9\)&\texttt{03215; 467}\\
\(3,4\)&\texttt{1740563; 2}&\(7,8\)&\texttt{5341027; 6}\\
\(3,5\)&\texttt{163; 25074}&\(7,9\)&\texttt{015; 64327}\\
\(3,6\)&\texttt{12; 476305}&\(8,9\)&\texttt{0235146; 7}\\
\bottomrule
\end{tabular}
\caption{Certificates for the two-color subgraphs.
Each word lists the vertices of a path in order; a semicolon
separates components. A one-letter word denotes an isolated vertex.}
\label{tab:paths}
\end{table}

\section{Nonextension}

For a vertex \(v\), let \(M(v)\) be the set of colors absent from
its incident edges. The matrix gives
\[
\begin{array}{c|cccccccc}
v&0&1&2&3&4&5&6&7\\ \hline
M(v)&\{2,9\}&\{2,3\}&\{3,4\}&\{4,5\}&
\{5,6\}&\{6,7\}&\{7,8\}&\{8,9\}.
\end{array}
\]

\begin{lemma}\label{lem:extensions}
After adding a vertex \(x\), the coloring \(c\) has exactly two
proper extensions with palette \(\{1,\ldots,9\}\).
Their respective lists of colors on \(x0,\ldots,x7\) are
\[
(2,3,4,5,6,7,8,9)
\quad\text{and}\quad
(9,2,3,4,5,6,7,8).
\]
\end{lemma}

\begin{proof}
Write \(f(v)\) for the color assigned to \(xv\).
Properness at \(v\) requires \(f(v)\in M(v)\), and properness at
\(x\) requires the eight values of \(f\) to be distinct.

If \(f(0)=2\), distinctness forces \(f(1)=3\), then \(f(2)=4\),
and successively \(f(3)=5,\ldots,f(7)=9\).
If \(f(0)=9\), distinctness instead forces \(f(7)=8\), then
\(f(6)=7\), and successively \(f(5)=6,\ldots,f(1)=2\).
These exhaust the two possibilities for \(f(0)\).
Each resulting assignment uses distinct colors and chooses a
member of \(M(v)\) at every vertex, so each is a proper extension.
\end{proof}

\begin{proof}[Proof of Theorem~\ref{thm:main}]
Lemma~\ref{lem:acyclic} establishes properness and acyclicity
of \(c\). By Lemma~\ref{lem:extensions}, only two proper extensions
need be considered.

In the first extension, \(f(0)=2\) and \(f(2)=4\).
The cycle
\[
x-0-4-5-6-2-x
\]
has successive edge colors \(2,4,2,4,2,4\).

In the second extension, \(f(1)=2\) and \(f(3)=4\).
The cycle
\[
x-1-7-3-x
\]
has successive edge colors \(2,4,2,4\).
Neither extension is acyclic.
\end{proof}

\begin{thebibliography}{1}

\bibitem{ASZ}
N.~Alon, B.~Sudakov, and A.~Zaks,
Acyclic edge colorings of graphs,
\emph{Journal of Graph Theory} \textbf{37} (2001), 157--167.
\href{https://doi.org/10.1002/jgt.1010}{doi:10.1002/jgt.1010}.

\end{thebibliography}
\end{document}
